\documentclass{article}
\usepackage[T1]{fontenc}
\usepackage{lmodern}
\usepackage{graphicx}
\usepackage{amsmath,amssymb}
\usepackage[a4paper,margin=2cm]{geometry}
\usepackage[absolute,overlay]{textpos}
\setlength{\TPHorizModule}{1mm}
\setlength{\TPVertModule}{1mm}
\usepackage[indonesian]{babel}
\usepackage{hyperref}
\setlength{\emergencystretch}{2em}
% unit: Halaman 33

\begin{document}

% === Halaman 33 (llm) ===

Dekripsi: Bagi panjang cipherteks dengan panjang kunci (Pada contoh ini, $30 / 6 = 5$).\\
Cipherteks: \textcolor{red}{SDNIAIAONSSNLFITTOOIEEGRTMKIMB}\
Tulis cipherteks secara vertical sepanjang 5 kolom:\
\vspace{5mm}
\begin{verbatim}
SDNIA
IAONS
SNLFI
TTOOI
EEGRT
MKIMB
\end{verbatim}
\vspace{5mm}
Plainteks: (baca secara vertikal)\
\texttt{sistemdanteknologiinformasiitb}\
\texttt{sistem dan teknologi informasi itb} \quad (tambahkan spasi)

\end{document}
